\subsection{Implementasi Nyata}
Algoritma String Matching seperti Boyer-Moore diterapkan secara nyata dalam berbagai aplikasi seperti pencarian teks di dokumen, mesin pencari, deteksi pola pada log keamanan, dan scanning barcode. Makalah berikut mendukung relevansi penerapan Boyer-Moore dalam aplikasi Barcode Generator untuk serial number matching.

\vspace{1em}
\textbf{[1] Penulis:} Tarigan, D. A., Buaton, A. O., Briyandana, Safitri, E. R., dan Rosnelly, R. (2024)

\textbf{Judul:} "Analysis of String Matching Application on Serial Number Using Boyer Moore Algorithm"

\textbf{Publikasi:} Journal of Computer Networks, Architecture and High Performance Computing, Vol. 6, No. 1, Hal. 237-246. DOI: 10.47709/cnahpc.v6i1.3410

\textbf{Intisari:} Paper ini membahas penerapan Boyer-Moore algorithm dalam aplikasi Barcode Generator untuk scanning barcode product serial numbers di Coway International Indonesia. Boyer-Moore algorithm digunakan untuk menemukan sequence numbers dari barcode yang telah diterjemahkan menjadi text. Algoritma ini dipilih karena lebih efisien dibanding binary search dan sequential search.

Hasil analisis menunjukkan string matching menggunakan Boyer-Moore memerlukan 9 iterasi untuk mencapai kondisi matched. Pattern dapat digeser hingga 13 kali dalam proses pencarian, menunjukkan efisiensi skip logic yang menjadi kekuatan Boyer-Moore.

\textbf{Link:} \url{https://doi.org/10.47709/cnahpc.v6i1.3410}

\vspace{1em}
\textbf{Relevansi dengan Praktikum:} Seperti pada praktikum String Matching, Boyer-Moore melakukan pencocokan dari kanan ke kiri dan menggunakan tabel last occurrence untuk menentukan pergeseran pattern. Kemampuan skip area yang tidak match sangat menguntungkan untuk serial number yang panjang.

\vspace{1em}
\textbf{Kelebihan dalam konteks barcode serial number scanning:}
\begin{itemize}
    \item Lebih efisien dibanding sequential search karena dapat melakukan skip berdasarkan last occurrence table, menghemat waktu scanning.
    \item Pada teks panjang seperti serial number, pergeseran pattern dapat lebih jauh sehingga mengurangi jumlah perbandingan karakter.
    \item Cocok untuk aplikasi real-time seperti barcode scanning yang membutuhkan response time cepat.
    \item Implementasi relatif sederhana dengan preprocessing last occurrence table yang straightforward.
\end{itemize}

\textbf{Kekurangan dalam konteks barcode serial number scanning:}
\begin{itemize}
    \item Untuk pattern yang sangat pendek (beberapa karakter), overhead preprocessing last occurrence table dapat mengurangi benefit.
    \item Worst case complexity tetap O(mn) meski average case sangat baik, sehingga pada kondisi tertentu performa bisa menurun.
    \item Jika terjadi kesalahan scanning barcode (noise), algoritma harus dijalankan ulang dari awal untuk pattern matching.
    \item Membutuhkan memory untuk menyimpan last occurrence table, meski ukurannya relatif kecil untuk alphabet yang terbatas.
\end{itemize}
